\documentclass[aps,prb,twocolumn,floatfix,nofootinbib,superscriptaddress]{revtex4-2}

\usepackage{amsmath,amssymb,mathtools,bm}
\usepackage{booktabs}
\usepackage[colorlinks=true,linkcolor=blue,citecolor=blue,urlcolor=blue]{hyperref}
\usepackage{microtype}

\newcommand{\tr}{\operatorname{tr}}
\newcommand{\ii}{\mathrm{i}}
\newcommand{\dd}{\mathrm{d}}
\newcommand{\code}[1]{\texttt{#1}}

\hypersetup{
  pdftitle={Adaptive Kato Transport of Overlap-Continued Circuit Projectors},
  pdfauthor={Asadullah Bhuiyan}
}

\begin{document}

\title{Adaptive Kato Transport of Overlap-Continued Circuit Projectors:\
A Preregistered Width Comparison}
\author{Asadullah Bhuiyan}
\affiliation{Class-A fermionic adaptive-circuit project}
\date{September 6, 2026}

\begin{abstract}
We specify a geometric flux-continuation experiment on saved pure Gaussian
states produced by the monitored class-A circuit.  At each flux, a fixed-rank
spectral subspace of the state-derived flattened parent is selected by maximum
overlap with the preceding projector.  Its frame is then transported with the
Kato connection rather than a physical-time Schr\"odinger ramp.  Adaptive
classical RK4 uses projector step doubling, while the selected--excluded
spectral gap is retained as a diagnostic.  The campaign compares
$20\times24$ and $24\times24$ endpoint ensembles, separately reporting raw
clockwise and counterclockwise wall-charge transfer.  Numerical validity is
gated, but charge quantization and agreement with earlier finite-grid
continuations are measured outcomes.  This document fixes the theory,
algorithm, diagnostics, and resume contract before production results are
inspected.
\end{abstract}

\maketitle

\section{Question and scope}

The equilibrium Laughlin argument identifies a Chern phase by transporting an
occupied spectral subspace around one boundary-flux quantum.  A monitored
circuit supplies several inequivalent flux-dependent objects.  A physical
online circuit ramp includes Born noise and feedback, and a finite-time
Schr\"odinger ramp of a state-derived Hamiltonian includes Landau--Zener
mixing.  Neither is identical to geometric continuation of a spectral branch.
The present experiment isolates the latter operation.

The input is a completed monitored endpoint state.  No circuit trajectory is
rerun, no tangent cocycle is constructed, and no physical ramp time is
introduced.  The output is the direct half-cylinder charge of a projector
transported continuously in flux.  This is the closest circuit-state analogue
of the overlap-continued flattened-Hamiltonian benchmark, but it does not make
quantization automatic: the endpoint family may lack a clean isolated branch,
the finite cylinder contains two walls, and previous-overlap selection can
become ambiguous near a degeneracy.

\section{Pure Gaussian states and their projectors}

A number-conserving pure Gaussian state with occupied frame
$F\in\mathbb{C}^{N\times r}$ satisfies
\begin{equation}
  F^\dagger F=\mathbf{1}_r,
  \qquad P=FF^\dagger,
  \qquad P^2=P=P^\dagger.
  \label{eq:frame-projector}
\end{equation}
The frame basis is not physical: $F$ and $FU$ represent the same state for
every $U\in U(r)$.  The rank-$r$ projector $P$ is the physical one-particle
correlation matrix.  Regional particle numbers therefore follow directly from
its diagonal,
\begin{equation}
  N_A=\tr(R_A P),
\end{equation}
where $R_A$ projects onto a spatial region.

For an endpoint projector $P_0$, define the state-derived flattened parent
\begin{equation}
  H_0=\mathbf{1}_N-2P_0.
  \label{eq:parent}
\end{equation}
It has occupied eigenvalue $-1$ and empty eigenvalue $+1$ at zero twist.
Threading a uniform boundary flux in the periodic $y$ direction gives
\begin{align}
  \widetilde H_{rs}(\phi)
    &=(H_0)_{rs}\exp\!\left[\ii\phi\frac{d_y(r,s)}{N_y}\right],\\
  H(\phi)&=\frac{\widetilde H(\phi)+\widetilde H^\dagger(\phi)}{2},
  \label{eq:twist}
\end{align}
with minimum-image displacement $d_y$.  The derivative is analytic,
\begin{align}
  \partial_\phi\widetilde H_{rs}
  &=\ii\frac{d_y(r,s)}{N_y}\widetilde H_{rs},\\
  \partial_\phi H
  &=\frac{\partial_\phi\widetilde H+
  (\partial_\phi\widetilde H)^\dagger}{2}.
  \label{eq:analytic-h-derivative}
\end{align}

\section{Three distinct flux-dependent projectors}

The notation distinguishes three operations that can otherwise look
deceptively similar.

The \emph{dynamical projector} $P_{\mathrm{dyn}}(t)$ solves a physical-time
equation such as
\begin{equation}
  \ii\partial_tF_{\mathrm{dyn}}=H[\phi(t)]F_{\mathrm{dyn}}.
\end{equation}
Its result depends on ramp time and can mix instantaneous branches at an
avoided crossing.

The \emph{instantaneous projector} $P_{\mathrm{inst}}(\phi)$ is rebuilt
independently at every flux by filling the $r$ lowest eigenvectors of
$H(\phi)$.  It has no memory of how the flux point was reached.  In a finite
system it normally closes after $2\pi$ by refilling according to energy order.

The \emph{continued projector} $P_{\mathrm{cont}}(\phi)$ remembers the
previous spectral subspace.  At a numerical stage with eigenpairs
$\{E_n,|w_n\rangle\}$, every eigenvector receives the gauge-invariant weight
\begin{equation}
  \omega_n=\langle w_n|P_{\mathrm{previous}}|w_n\rangle.
  \label{eq:weights}
\end{equation}
The fixed-rank set $I_{\mathrm{cont}}$ consists of the $r$ largest weights.
Only exact numerical ties are resolved by increasing energy, reproducing the
legacy static selector.  The selected spectral projector is
\begin{equation}
  P_{\mathrm{spec}}(\phi)=
  \sum_{i\in I_{\mathrm{cont}}}|i\rangle\langle i|.
  \label{eq:pspec}
\end{equation}
If the branch is smooth and resolved, Kato transport makes
$P_{\mathrm{cont}}=P_{\mathrm{spec}}$ up to integration error.  We keep the
two names because their discrepancy is a primary numerical diagnostic.

\section{Kato parallel transport}

\subsection{Derivation from $P^2=P$}

Differentiating $P^2=P$ gives
\begin{equation}
  (\partial_\phi P)P+P(\partial_\phi P)=\partial_\phi P,
  \qquad P(\partial_\phi P)P=0.
  \label{eq:projector-identity}
\end{equation}
Thus $\partial_\phi P$ is off diagonal between the selected subspace and its
complement.  Define the anti-Hermitian Kato generator
\begin{equation}
  K(\phi)=[\partial_\phi P_{\mathrm{spec}},P_{\mathrm{spec}}],
  \qquad K^\dagger=-K.
  \label{eq:kato-generator}
\end{equation}
The frame equation is
\begin{equation}
  \partial_\phi F_{\mathrm{cont}}=K(\phi)F_{\mathrm{cont}}.
  \label{eq:frame-kato}
\end{equation}
It induces the basis-independent projector equation
\begin{equation}
  \partial_\phi P_{\mathrm{cont}}
  =[[\partial_\phi P_{\mathrm{spec}},P_{\mathrm{spec}}],
  P_{\mathrm{cont}}].
  \label{eq:projector-kato}
\end{equation}
Indeed, $\partial_\phi(FF^\dagger)=KP-PK=[K,P]$; when
$P=P_{\mathrm{spec}}$, Eq.~\eqref{eq:projector-identity} gives
$[K,P]=\partial_\phi P$.  The connection fixes a convenient parallel gauge
for the frame but transports the projector, which is the physical object.

\subsection{Analytic spectral derivative}

Away from an unresolved selected--excluded degeneracy, first-order spectral
perturbation theory gives
\begin{equation}
  \partial_\phi P_{\mathrm{spec}}
  =\sum_{\substack{i\in I_{\mathrm{cont}}\\a\notin I_{\mathrm{cont}}}}
  \frac{|a\rangle\langle a|(\partial_\phi H)|i\rangle\langle i|}
  {E_i-E_a}+\mathrm{H.c.}
  \label{eq:spectral-derivative}
\end{equation}
The selected--excluded gap
\begin{equation}
  \delta_{\mathrm{sel}}=
  \min_{\substack{i\in I_{\mathrm{cont}}\\a\notin I_{\mathrm{cont}}}}
  |E_i-E_a|
  \label{eq:selected-gap}
\end{equation}
controls the norm that Eq.~\eqref{eq:spectral-derivative} can acquire.  It is
not, by itself, an integration-error estimate.  Matrix elements of
$\partial_\phi H$, rotation distributed among many modes, and branch changes
also matter.  For that reason $\delta_{\mathrm{sel}}$ is recorded but does not
set the flux step directly.

\section{Adaptive RK4 algorithm}

The regulated observation grid is
\begin{equation}
  \phi_j^{(\sigma)}=-\sigma10^{-7}
  +\sigma\frac{2\pi j}{128},
  \quad j=0,\ldots,128,
  \quad \sigma=\pm1.
  \label{eq:grid}
\end{equation}
The signed $10^{-7}$ displacement starts on a definite side of the ambiguous
finite-size crossing at the nominal origin.  It does not alter the Hamiltonian
away from that negligible shift and is not a quantization criterion.

For each RK4 stage, the code constructs $H$ and its analytic derivative,
diagonalizes $H$, performs the previous-projector selection, constructs
Eq.~\eqref{eq:spectral-derivative}, and applies the generator in
Eq.~\eqref{eq:frame-kato}.  The same proposed step is evaluated once at size
$h$ and twice at size $h/2$.  After symmetric polar orthonormalization of the
two candidates, the local projector discrepancy is
\begin{equation}
  e=\frac{\|P_{h/2,h/2}-P_h\|_F}{\sqrt r}.
  \label{eq:step-error}
\end{equation}
A step is accepted only for $e\le10^{-8}$.  The next proposal is
\begin{equation}
  h_{\mathrm{new}}=0.9h\left(\frac{10^{-8}}{e}\right)^{1/5},
  \qquad 0.2h\le h_{\mathrm{new}}\le2h.
  \label{eq:controller}
\end{equation}
No step may cross an observation point.  Failure is diagnostic if the
proposal falls below $1/4096$ of one observation interval.  This rule catches
an unresolved branch or singular derivative instead of silently stepping
through it.

After every accepted step, the frame is stabilized by the symmetric polar
map
\begin{equation}
  F\longmapsto F(F^\dagger F)^{-1/2}.
  \label{eq:polar}
\end{equation}
The result is \emph{not} reprojected onto $P_{\mathrm{spec}}$.  Reprojection
would hide accumulated transport error and could convert a branch-selection
discontinuity into an apparently exact path.  The direct projector mismatch
$\|P_{\mathrm{cont}}-P_{\mathrm{spec}}\|_F/\sqrt r$ remains observable and
must stay below $10^{-7}$.

\section{Campaign and saved observables}

The inputs are verified complex128 endpoint frames from two completed raster-
$y$ monitored campaigns.  Both used pure initialization, perfect correction,
no postselection, $n_{\mathrm{shell}}=1$, and
$(\alpha_1,\alpha_2)=(1,30)$ for 48 cycles.  The width comparison is
\begin{center}
\begin{tabular}{lccc}
\toprule
size & walls & split & selected IDs\\
\midrule
$20\times24$ & $x=5,15$ & $x=10$ & $0,4,\ldots,96$\\
$24\times24$ & $x=6,18$ & $x=12$ & $0,4,\ldots,96$\\
\bottomrule
\end{tabular}
\end{center}
Soft and hard walls are both included.  Each of the 100 endpoint states is
continued clockwise and counterclockwise, yielding 200 paths.  The two sizes
are independent monitored ensembles; the two directions for a fixed endpoint
are paired deterministic continuations.

At all 129 points we save $\phi$, $N_L$, $N_R$, total charge,
\begin{equation}
  q_x=\frac{(N_R-N_R^{(0)})-(N_L-N_L^{(0)})}{2},
\end{equation}
the $x$-resolved density, transported--selected projector mismatch, minimum
principal overlap, selected-weight floor and margin, $\delta_{\mathrm{sel}}$,
generator norm, adaptive error, and accepted/rejected step counts.  Raw CW
and CCW distributions are reported separately.  The analysis compares each
path to its saved 64-interval static continuation and also to the 128-interval
static continuation at $20\times24$.  Every change in the descriptive
$|q_x|>0.5$ classification is listed.  Fractions above $0.5$ and $0.75$ are
descriptive outcomes, not gates.

Each path writes one atomic compressed NPZ and one completion JSON written
last.  The completion binds the task, full source endpoint identity,
configuration hash, source hashes, result bytes, and SHA-256.  A rolling
uncompressed frame checkpoint is written every eight observation intervals
as one stable NPZ/JSON pair.  It contains the exact transported frame,
completed interval, next adaptive proposal, accumulated diagnostics, and all
partial histories.  Restart verifies and skips complete pairs or continues a
matching checkpoint; corrupt or partial pairs are rejected.

\section{Numerical gates and interpretation}

The preregistered numerical gates are: source-frame Gram residual below
$10^{-8}$; post-polar orthonormality residual below $10^{-10}$; total-charge
residual below $10^{-9}$; transported--selected projector mismatch below
$10^{-7}$; finite spectral and adaptive diagnostics; and step-doubling error
not exceeding $10^{-8}$.  Sample 0 for both sizes, walls, and directions is
repeated at tolerance $2.5\times10^{-9}$, with full-path and endpoint raw
$q_x$ differences required below $10^{-4}$.

Kato continuation removes physical ramp time and hence removes finite-rate
Landau--Zener mixing from the definition of the path.  It cannot force charge
quantization.  A true ambiguity can remain if the selected--excluded gap
closes, if the overlap margin vanishes, or if a finite-system continued branch
does not carry a full unit of wall charge.  Quantization, agreement with the
coarse-grid classification, and width improvement are therefore reported
only after all 200 path pairs verify.  No result claim is made in this methods
note.

\begin{thebibliography}{9}
\bibitem{Kato1950}
T. Kato, On the adiabatic theorem of quantum mechanics,
J. Phys. Soc. Jpn. \textbf{5}, 435 (1950).
\bibitem{Laughlin1981}
R. B. Laughlin, Quantized Hall conductivity in two dimensions,
Phys. Rev. B \textbf{23}, 5632 (1981).
\bibitem{BhuiyanPanJian2026}
A. Bhuiyan, L. Pan, and C.-M. Jian, Free-fermion dynamics with measurements:
Topological classification and adaptive preparation of topological states,
Phys. Rev. Research \textbf{8}, 023147 (2026).
\end{thebibliography}

\end{document}
